\documentclass{ifpr-pinhais}

\usepackage{amsmath}
\usepackage{graphicx}
\usepackage{latexsym}
\usepackage{amssymb}
\usepackage{placeins}
\usepackage{listings}

\graphicspath{{imagens/}}
\lstset{basicstyle=\ttfamily\footnotesize,columns=fixed,keepspaces=true,upquote=true,breaklines=true,postbreak=\mbox{$\hookrightarrow$\space},frame=single,showstringspaces=false}
\newcolumntype{P}[1]{>{\raggedright\arraybackslash}p{#1}}

% ---- Metadados do trabalho ----
\curso{Bacharelado em Ciência da Computação}
\departamento{Coordenação do Curso de Bacharelado em Ciência da Computação}
\autor{Beatriz Regina Nahas \\ João Pedro dos Santos Henrique Plinta}
\titulo{Manual 1 -- eBPF para Observabilidade}
\tipotrabalho{Manual}
\local{Pinhais}
\data{2026}

\orientador{Prof. Guilherme Werneck de Oliveira}

\preambulo{%
  Manual produzido no âmbito do projeto de extensão do minicurso
  Introdução ao eBPF, apresentado à disciplina de Práticas de Extensão V,
  do curso de Bacharelado em Ciência da Computação do \ifpr, \campus,
  sob a orientação do professor Guilherme Werneck de Oliveira.
}

\makeatletter
\hypersetup{
  pdftitle={\@title},
  pdfauthor={\@author},
  pdfsubject={\imprimirpreambulo},
  pdfcreator={LaTeX com abntex2 e ifpr-pinhais},
  hidelinks,
}
\makeatother

\setlength{\parskip}{0.1cm}

\makeindex

\begin{document}
\frenchspacing

\pretextual
\imprimircapa
\imprimirfolhaderosto{\folhaderosto}
\pdfbookmark[0]{\contentsname}{toc}
\tableofcontents*
\cleardoublepage

\textual
\pagestyle{simple}

\chapter{Introdução}

\section{O que é eBPF}

\section{Arquitetura e funcionamento}

\section{eBPF aplicado à observabilidade}

\chapter{Pré-requisitos}

\section{Conhecimentos necessários}

\section{Ferramentas necessárias}

\section{Instalação e configuração do ambiente}

\chapter{Fundamentos de Observabilidade}

\section{O que é observabilidade}

\section{Métricas}

\section{Traces}

\section{Eventos}

\chapter{eBPF para Monitoramento}

\section{Rastreamento de processos e chamadas de sistema}

\section{Monitoramento de CPU e memória}

\section{Monitoramento de disco e I/O}

\section{Monitoramento de aplicações}

\chapter{Ferramentas de Observabilidade}

\section{BCC}

\section{bpftrace}

\section{libbpf}

\section{eBPF Exporters}

\chapter{Instrumentação com eBPF}

\section{Kprobes}

\section{Kretprobes}

\section{Tracepoints}

\section{Uprobes}

\section{USDT}

\section{Perf Events}

\chapter{Coleta e Visualização de Dados}

\section{eBPF e Prometheus}

\section{eBPF e Grafana}

\section{Exportação de métricas}

\section{Coleta de eventos}

\section{Dashboards de observabilidade}

\chapter{Exemplos Práticos}\label{cap:exemplos}

Este capítulo reúne três cenários que combinam as ferramentas dos capítulos anteriores. Cada um parte de uma pergunta de diagnóstico e usa só comandos já apresentados, de modo que o leitor possa reproduzi-los em um terminal com dois painéis: um para a ferramenta de coleta e outro para gerar a atividade observada. Todos exigem \texttt{root} e foram testados no ambiente descrito no capítulo \ref{cap:monitoramento}.

\section{Monitoramento de processos}

\textbf{Pergunta:} que processos de vida curta estão sendo criados no sistema? Um \textit{script} que roda por milissegundos, disparado pelo \texttt{cron} ou por um serviço, não aparece em um \texttt{ps} nem em uma amostra de \textit{polling}.

No primeiro terminal, inicie o rastreamento de \texttt{execve} e deixe-o rodando:

\begin{lstlisting}
sudo execsnoop-libbpf
\end{lstlisting}

No segundo terminal, crie alguns processos curtos:

\begin{lstlisting}
for i in 1 2 3; do date > /dev/null; sleep 0.2; done
\end{lstlisting}

O primeiro terminal passa a listar cada \texttt{date} e cada \texttt{sleep}, com o PID, o processo pai e os argumentos. O \texttt{ps} não mostraria nenhum deles. Para saber \textit{quem} está criando processos em maior volume, em vez de listar cada um, troque a ferramenta por uma contagem:

\begin{lstlisting}
sudo bpftrace -e 'tracepoint:syscalls:sys_enter_execve { @[comm] = count(); }'
\end{lstlisting}

Ao encerrar com \texttt{Ctrl+C}, o \texttt{bpftrace} imprime o mapa com o total de execuções por processo pai. Um número alto em um \textit{script} de verificação periódica é um indício de que ele cria processos demais para o que faz.

\section{Análise de chamadas de sistema}

\textbf{Pergunta:} o que um programa está pedindo ao kernel? Aqui o alvo é o servidor HTTP embutido do Python. Inicie-o no primeiro terminal:

\begin{lstlisting}
python3 -m http.server 8080
\end{lstlisting}

No segundo terminal, descubra o PID do servidor e observe os arquivos que ele abre enquanto atende requisições:

\begin{lstlisting}
PID=$(pgrep -f "http.server 8080")
sudo opensnoop-libbpf -p $PID
\end{lstlisting}

Em um terceiro terminal, faça requisições com \texttt{curl http://localhost:8080/}. Cada requisição aparece no \texttt{opensnoop} como a abertura do arquivo ou do diretório pedido, e é possível ver o que o servidor lê do disco para responder.

Para a visão agregada, o \texttt{syscount} mostra quais chamadas o servidor mais faz, sem interromper o processo, o que o \texttt{strace} não consegue evitar:

\begin{lstlisting}
sudo syscount-libbpf -p $PID -i 5
\end{lstlisting}

Comparar essa saída antes e depois de uma carga de requisições revela a proporção entre chamadas de rede (\texttt{sendto}, \texttt{recvfrom}), de arquivo (\texttt{openat}, \texttt{read}) e de espera (\texttt{epoll\_wait}, \texttt{poll}). Uma espera dominante indica um servidor ocioso, não lento.

\section{Rastreamento de uma aplicação}

\textbf{Pergunta:} quem se conecta ao serviço, por quanto tempo e com quantos bytes? Com o mesmo servidor na porta 8080, rode o \texttt{tcplife}, que registra uma linha ao fim de cada conexão:

\begin{lstlisting}
sudo tcplife-libbpf -L 8080
\end{lstlisting}

Gere conexões em outro terminal:

\begin{lstlisting}
for i in $(seq 1 5); do curl -s http://localhost:8080/ > /dev/null; done
\end{lstlisting}

Cada linha traz o processo, os endereços e portas, os bytes enviados e recebidos e a duração. Conexões longas ou com muitos bytes se destacam do conjunto. Para ver também quem \textit{abre} conexões a partir da máquina, e não só as recebidas, use o \texttt{tcpconnect-libbpf} em paralelo.

Esses comandos observam o serviço de fora, pelo que o kernel vê. Para olhar por dentro de um processo Python, sem alterá-lo, use a sonda USDT do interpretador (seção \ref{sec:uprobes} e capítulo \ref{cap:instrumentacao}). O comando a seguir lista os módulos importados por qualquer programa Python, e ajuda a explicar um início lento:

\begin{lstlisting}
sudo bpftrace -e 'usdt:/usr/bin/python3.14:python:import__find__load__start {
    printf("%s\n", str(arg0)); }'
\end{lstlisting}

Em outro terminal, execute \texttt{python3 -c "import json, csv"} e observe os módulos aparecerem na ordem em que o interpretador os carrega. Para transformar uma dessas medidas em uma série contínua em um painel, siga o caminho do capítulo \ref{cap:coleta}.

\chapter{Troubleshooting}\label{cap:troubleshooting}

Este capítulo reúne os problemas mais prováveis ao seguir o manual e os limites da técnica. Os erros descritos foram observados nos testes em Ubuntu 26.04 com kernel 7.0.

\section{Impacto de performance}

Um programa eBPF roda toda vez que o evento ao qual está anexado ocorre, então o custo depende da frequência do evento e do trabalho feito a cada execução. Algumas regras práticas:

\begin{itemize}
  \item \textbf{Prefira eventos raros.} Um \textit{tracepoint} em \texttt{execve} dispara algumas vezes por segundo. Um em \texttt{raw\_syscalls:sys\_enter} dispara a cada chamada de sistema do host inteiro, o que pode chegar a milhões por segundo.
  \item \textbf{Agregue no kernel.} Contar com \texttt{count()} ou \texttt{hist()} custa muito menos do que imprimir cada evento com \texttt{printf}, porque o espaço do usuário só lê o resultado ao final.
  \item \textbf{Prefira amostragem a rastrear tudo.} Para saber onde a CPU é gasta, \texttt{profile:hz:99} tem custo fixo, independente da carga.
  \item \textbf{Cuidado com ferramentas pesadas.} O \texttt{memleak} instrumenta cada alocação do processo, e o \texttt{kretprobe} custa mais que o \texttt{kprobe}. Use-os por períodos curtos e, se possível, fora de produção.
\end{itemize}

Para medir o custo de uma sonda, rode a carga de trabalho com e sem ela e compare a latência ou a vazão. O resultado do experimento citado no capítulo \ref{cap:introducao} vale para o coletor testado, não para qualquer programa.

\section{Limitações}

\textbf{Permissões.} Carregar programas eBPF exige \texttt{root} (ou as capacidades \texttt{CAP\_BPF} e \texttt{CAP\_PERFMON}). Sem isso, as ferramentas falham com erro de permissão. Em contêineres, o acesso costuma ser bloqueado por padrão.

\textbf{BCC e cabeçalhos do kernel.} O BCC compila o programa na hora da execução, a partir dos cabeçalhos do kernel em uso. No kernel 7.0 do ambiente de teste, 11 das 14 ferramentas testadas do BCC falharam ao compilar, com erros como \texttt{static assertion failed} em \texttt{include/linux/fs.h}. As versões \texttt{-libbpf}, que usam CO-RE, funcionaram em 13 das 14. Se uma ferramenta do BCC não compilar, procure a equivalente com o sufixo \texttt{-libbpf}.

\textbf{Argumentos diferentes entre versões.} A mesma ferramenta pode aceitar argumentos distintos na versão BCC e na \texttt{libbpf}. O \texttt{ext4slower-libbpf 10}, por exemplo, falhou com \texttt{Too many arguments}. Consulte sempre \texttt{--help} da versão instalada.

\textbf{Avisos do bpftrace.} Ao usar USDT, o \texttt{bpftrace} 0.25.0 imprime avisos de compilação vindos do \texttt{vmlinux.h} (\texttt{declaration does not declare anything}). São inofensivos e a opção \texttt{--no-warnings} não os suprime. Já um programa USDT de várias linhas, salvo em arquivo, fez o \texttt{bpftrace} abortar com \texttt{Assertion failed}. Escreva \textit{scripts} USDT em uma única linha com \texttt{-e}.

\textbf{Sondas que não existem.} Uma sonda pode estar ausente porque o pacote foi compilado sem ela. O erro típico é \texttt{No matches for usdt}. Liste as sondas com \texttt{bpftrace -l} antes de escrever o programa.

\textbf{Funções internas mudam.} Um \textit{kprobe} depende do nome e dos argumentos de uma função interna do kernel, que não são uma interface estável. Se um programa deixar de funcionar após uma atualização, verifique primeiro se a função ainda existe com \texttt{bpftrace -l 'kprobe:nome*'}.

\textbf{Contadores de hardware.} Eventos como falhas de cache dependem do processador e geralmente não estão disponíveis em máquinas virtuais.

\textbf{O que o eBPF não vê.} A coleta enxerga o que passa pelo kernel. Eventos internos de uma aplicação, como o tempo de uma consulta SQL, exigem sondas \texttt{uprobe} ou USDT, ou instrumentação no próprio código.

\chapter{Conclusão}\label{cap:conclusao}

Este manual percorreu o caminho de uma observação ao primeiro painel de observabilidade com eBPF. Partiu dos conceitos (capítulos \ref{cap:introducao} e \ref{cap:monitoramento}), passou pelas ferramentas e pelos mecanismos de instrumentação (capítulos \ref{cap:ferramentas} e \ref{cap:instrumentacao}) e terminou na exportação de métricas para o Prometheus e na visualização no Grafana (capítulo \ref{cap:coleta}). Os exemplos práticos do capítulo \ref{cap:exemplos} mostram que as ferramentas prontas, como \texttt{execsnoop}, \texttt{opensnoop}, \texttt{syscount} e \texttt{tcplife}, respondem à maior parte das perguntas de diagnóstico sem escrever código.

O eBPF não substitui as ferramentas tradicionais. Ele complementa o \texttt{top}, o \texttt{iostat} e o Prometheus onde eles são mais fracos: eventos de curta duração, distribuições de latência em vez de médias e métricas por processo ou por conexão, coletadas sem alterar a aplicação. Em troca, exige \texttt{root}, um kernel recente e atenção ao custo de cada sonda, como o capítulo \ref{cap:troubleshooting} detalha.

Os testes que sustentam os comandos deste manual estão na pasta \texttt{testes/} do repositório, com a saída de cada execução. Como os resultados dependem do kernel e das versões das ferramentas, o leitor deve repeti-los em seu próprio ambiente antes de tomar qualquer número como referência. Como próximos passos, os Manuais 2 e 3 desta série aplicam o mesmo conjunto de técnicas a segurança e a redes.


\postextual

\bibliography{referencias}

\end{document}
